% =====================================================================
%  CHƯƠNG 1 — GIỚI THIỆU ĐỀ TÀI
% ---------------------------------------------------------------------
%  Người phụ trách : [TODO: Họ tên - MSSV]
%  Tiền tố label   : c1      (fig:c1:..., tab:c1:..., eq:c1:..., sec:c1:...)
%  Thư mục ảnh     : image/chuong1/
%  File TLTK       : references/chuong1.bib   (key: c1_...)
%  Độ dài gợi ý    : 2–4 trang
%  Mục tiêu        : Nêu bối cảnh, yêu cầu đề tài, thông số mục tiêu và
%                    sơ đồ khối tổng quát của hệ thống.
% =====================================================================
\chapter{Giới thiệu đề tài}
\label{chap:c1}

% ---------------------------------------------------------------------
\section{Đặt vấn đề}
\label{sec:c1:dat-van-de}
% Gợi ý: nhu cầu đo khối lượng chính xác (cân tiểu ly), ứng dụng chuẩn
% truyền tín hiệu 4--20mA trong công nghiệp, lý do chọn giải pháp.

\TODO{Viết 1--2 đoạn nêu bối cảnh và ý nghĩa của đề tài.}

% ---------------------------------------------------------------------
\section{Yêu cầu đề tài}
\label{sec:c1:yeu-cau}

\noindent\textbf{Đề \SoDeTai:} Thiết kế mạch đo khối lượng (tiểu ly) sử dụng Load cell.
Mạch có tầm đo \qtyrange{0}{1}{\kilo\gram}, độ phân giải \qty{2}{\milli\gram},
ngõ ra là dòng điện tương ứng từ \qtyrange{4}{20}{\milli\ampere} (thả nổi):
\begin{enumerate}[label=\alph*.]
    \item Lựa chọn và mô phỏng Load cell.
    \item Lựa chọn OPAMP sử dụng, mô phỏng OPAMP theo datasheet.
    \item Thiết kế mạch bù nhiệt độ.
    \item Đảm bảo sai số (giữa giá trị lý thuyết và thực tế đo qua mạch) $\pm\qty{5}{\percent}$.
    \item Thiết kế mạch đọc dòng điện và cho ngõ ra ADC. Tiến hành lựa chọn ADC
          (số bit, nguồn). Mô phỏng mạch đọc và kiểm chứng.
\end{enumerate}

% ---------------------------------------------------------------------
\section{Thông số kỹ thuật mục tiêu}
\label{sec:c1:thong-so}

Các thông số kỹ thuật mà hệ thống cần đạt được tổng hợp trong Bảng~\ref{tab:c1:thong-so-muc-tieu}.

\begin{table}[H]
    \centering
    \caption{Thông số kỹ thuật mục tiêu của hệ thống}
    \label{tab:c1:thong-so-muc-tieu}
    \renewcommand{\arraystretch}{1.3}
    \begin{tabular}{|l|c|}
        \hline
        \textbf{Thông số} & \textbf{Giá trị} \\ \hline
        Tầm đo                       & \qtyrange{0}{1}{\kilo\gram} \\ \hline
        Độ phân giải                 & \qty{2}{\milli\gram} \\ \hline
        Tín hiệu ngõ ra              & \qtyrange{4}{20}{\milli\ampere} (thả nổi) \\ \hline
        Sai số cho phép              & $\pm\qty{5}{\percent}$ \\ \hline
        Ngõ ra số                    & ADC \TODO{số bit} \\ \hline
    \end{tabular}
\end{table}

% ---------------------------------------------------------------------
\section{Sơ đồ khối tổng quát}
\label{sec:c1:so-do-khoi}
% Gợi ý: Load cell -> Khuếch đại vi sai -> Bù nhiệt -> V-I 4--20mA
%        -> Mạch đọc dòng -> ADC. Giải thích ngắn chức năng từng khối.

Sơ đồ khối tổng quát của hệ thống được trình bày trong Hình~\ref{fig:c1:so-do-khoi}.

\hinh[0.9\textwidth]{chuong1/so_do_khoi.png}{Sơ đồ khối tổng quát của hệ thống}{fig:c1:so-do-khoi}

\TODO{Mô tả chức năng của từng khối trong sơ đồ.}

% ---------------------------------------------------------------------
\section{Bố cục báo cáo}
\label{sec:c1:bo-cuc}

Báo cáo gồm các chương sau:
\begin{itemize}
    \item \textbf{Chương~\ref{chap:c1}} giới thiệu đề tài, yêu cầu và sơ đồ khối tổng quát.
    \item \textbf{Chương~\ref{chap:c2}} trình bày cơ sở lý thuyết và lựa chọn linh kiện.
    \item \textbf{Chương~\ref{chap:c3}} tính toán và mô phỏng các khối xử lý tín hiệu analog.
    \item \textbf{Chương~\ref{chap:c4}} thiết kế mạch đọc dòng và bộ chuyển đổi ADC.
    \item \textbf{Chương~\ref{chap:c5}} thiết kế và thi công phần cứng.
    \item \textbf{Chương~\ref{chap:c6}} trình bày kết quả đo thực tế và đánh giá sai số.
\end{itemize}
